% Appendix G -- the DirectKV ordering argument in full, moved out of section 3
% (2026-09-26). The body states the account and that we hold it weakly; this is
% why, in detail. The section itself says the contribution does not rest on it,
% which is the argument for it not occupying body pages.

\section{The DirectKV ordering, in detail}
\label{app:directkv}

DirectKV~\cite{directkv} reports that on Llama-3.1-8B its zero-copy host-KV
attention is \emph{slightly slower} than the HBM-resident SGLang baseline, and
that it wins on OPT-13B and OPT-30B, where SGLang runs out of memory.

Our amplifier offers an account of the ordering. Llama-3.1-8B's 16\,GB of
weights leave $\Ckv \approx 76$\,GB of a 96\,GB part, so
$\HBM/\Ckv \approx 1.26$ and there is little to buy; any overhead the mechanism
carries then shows up as a net loss. OPT-30B's 60\,GB leave
$\Ckv \approx 32$\,GB, tripling the amplifier.

We put this no more strongly than that, for three reasons a reader will find
anyway. First, Eq.~\ref{eq:gain} compares a \emph{mixture} against all-HBM,
while DirectKV is all-host; our law does not predict that an all-host system
beats an all-HBM one, and on its own terms it predicts the opposite. Second,
their OPT wins are reported ``where SGLang OOMs'', and OPT uses multi-head
attention with several times the KV bytes per token of a GQA model --- so the
baseline may be failing because $S\btok$ is large rather than because weights
are, and those are different axes. Third, OPT-13B has $\HBM/\Ckv$ close to the
8B's and yet DirectKV wins on it, which our amplifier alone does not account
for. The honest summary is that the model is \emph{consistent with} the
direction of a two-point published ordering, which is weak evidence, and we do
not rest the contribution on it.
